
\section{Related Work}
\label{sec:related}

\textbf{Play, Curiosity, and Developmental Robotics.}
Play has long been viewed as a central mechanism for skill development in natural intelligence: before children are given explicit tasks, they explore objects, discover controllable effects, and practice emerging motor routines~\citep{piaget1952origins, vygotsky1978mind, pellegrini2009role, gopnik2020childhood}. Rather than random activity, play provides a self-directed curriculum that samples interactions that are novel, meaningful, and near the boundary of current competence. This view has inspired computational models of curiosity and intrinsic motivation, where agents acquire skills before external rewards by pursuing prediction error, information gain, novelty, learning progress, or intermediate difficulty~\citep{schmidhuber1991curious, schmidhuber2010formal, houthooft2016vime, oudeyer2007intrinsic, kaplan2007intrinsic, kidd2012goldilocks}. Developmental robotics instantiated these ideas through Intelligent Adaptive Curiosity, goal babbling, competence-progress-driven goal selection, and modular curiosity systems for discovering object interactions and tool-use precursors~\citep{rolf2010goal, baranes2013active, forestier2016modular, forestier2017intrinsically}. In robot learning, play has provided broad, task-agnostic interaction data for learning reusable manipulation behaviors, latent plans, and hierarchical policies without segmented task demonstrations~\citep{lynch2020learning, wang2023mimicplay}, while language-conditioned curiosity uses compositional goals to expand exploration~\citep{colas2020language}. \textsc{RATs} extends this play-driven skill-acquisition paradigm to Code-as-Policy agents, where self-proposed practice goals are executed as robot programs and distilled into reusable code skills.

\textbf{Continual Skill Learning and Curriculum in Robotics.}
Continual robot learning studies how agents acquire, retain, and reuse knowledge over extended experience~\citep{thrun1995lifelong, parisi2019continual, lesort2020continual}. In manipulation, this often takes the form of reusable skills, options, motor primitives, affordances, behavioral priors, or hierarchical policies that can be transferred and composed across tasks~\citep{sutton1999between, konidaris2012robot, kroemer2021review, pertsch2020long, lynch2021language, wan2024lotuscontinualimitationlearning, ma2024dreureka}. Curriculum learning provides a complementary mechanism for improving long-horizon learning by ordering tasks from simple to complex~\citep{bengio2009curriculum, narvekar2020curriculum}, selecting tasks, goals, or demonstrations near the agent's competence boundary, expanding from known states, or training goal-conditioned policies over diverse goals~\citep{florensa2018automatic, nair2018visual, pong2020skewfit, colas2020language, seita2019zpd}. These methods provide mechanisms for skill reuse and automatic task ordering, but typically assume a predefined task family, goal representation, reward function, or externally provided experience stream. \textsc{RATs} differs by treating the curriculum itself as an object of autonomous discovery in a Code-as-Policy agent.


\textbf{Code-as-Policy and Agentic Robot Learning.}
Code-as-Policy methods use large language or multimodal models to synthesize executable robot programs that compose perception, planning, and control APIs, making robot policies modular, inspectable, and reusable~\citep{liang2023code, singh2023progprompt, vemprala2024chatgpt, mu2024robocodexmultimodalcodegeneration, geminirobotics1d52025, shi2025maestro}. Related embodied LMM systems ground language plans in pretrained skills, affordance models, scene representations, visual feedback, or 3D value maps for long-horizon manipulation~\citep{ahn2022saycan, driess2023palme, huang2023voxposer, zeng2023socratic, mu2023embodied}. Agentic methods further use reasoning-action loops, execution feedback, self-refinement, and tool use to improve generated programs~\citep{yao2023react, shinn2023reflexion, madaan2023selfrefine, chen2024selfdebug}; in robotics, CaP-X shows that multi-turn feedback, visual differencing, parallel reasoning, skill synthesis, and verifiable rewards improve embodied coding agents~\citep{fu2026capx}. However, most robot Code-as-Policy systems remain task-driven: an agent writes and revises code for a given instruction, and reusable skills are accumulated only as byproducts of solving provided tasks. Open-ended LMM agents such as Voyager show that automatic curricula and executable skill libraries can support lifelong exploration~\citep{wang2023voyager}, while sleep-time compute suggests spending computation before queries arrive to improve future performance~\citep{lin2025sleeptime}. \textsc{RATs} brings this skill-acquisition view to physically grounded Code-as-Policy learning, asking what the robot should practice before downstream tasks are specified.


